\paragraph{M1 for citations.}
Annual and cumulative citation networks share the same node set at year $t$, namely all
selected papers published by $t$. Directed density is $E/[N(N-1)]$ and mean total degree
is $2E/N$; clustering, component counts and LCC fractions use the weak undirected
projection, including isolates. Trailing windows of 1, 2, 5 and 10 years provide the
window grid. Window sensitivity, $(\max-\min)/\mathrm{mean}$ over the grid, is computed at
each end year and averaged over end years 1990--2020, which constitute the primary
analysis interval; earlier papers enter as history.

\paragraph{M2c-1: preferential attachment.}
Following the degree-class approach of \citet{jeong2003measuring}, in-degrees are frozen
at the start of each citing year, and for each prior-degree class $k$ the number of
eligible papers $n_{tk}$ and of new citations $h_{tk}$ are counted. The exponent $\alpha$
is estimated by conditional likelihood with target probabilities proportional to
$n_{tk}(k+1)^\alpha$, conditioning on the total number of citations in each year. The
offset admits citations to previously uncited papers; offsets of 0.5 and 2 serve as
sensitivity checks. Confidence intervals use 400 bootstrap resamples of citing-year blocks
and are descriptive. Neither a positive $\alpha$ nor its departure from an age-only
reference identifies a causal mechanism.

\paragraph{M2c-2: ageing.}
For ages 1--60, the target probability in citing year $t$ is proportional to
$n_{ta}f(a)$, where $n_{ta}$ is the number of available papers of age $a$. Lognormal,
exponential, power-law and uniform kernels $f$ are compared by AIC on identical
observations and support. The comparison is repeated by publication cohort (1970--1989,
1990--1999, 2000--2009), accounting for each cohort's exposure. The estimate describes
marginal ageing, which may also reflect degree, fitness and venue coverage.

\paragraph{M2c-3: densification and diameter.}
The densification exponent $a$ is obtained from $\log E(t)=a\log N(t)+b$ over 1990--2020
and over split periods and an earlier cutoff. Year-resampling intervals are descriptive
because cumulative observations are dependent. The effective diameter is the 90th
percentile of shortest-path distances within the largest weak component, estimated from
128 uniformly sampled source nodes per year, with 64 and 256 sources as sensitivity checks
\citep{leskovec2005graphs}.

\paragraph{M2c-4: citation dynamics and prediction.}
The Wang--Song--Barab\'asi (WSB) model \citep{wang2013quantifying},
\begin{equation}
 c(t)=m\left[\exp\!\left\{\lambda\Phi\left(\frac{\log t-\mu}{\sigma}\right)\right\}-1\right],
\end{equation}
is fitted with $m$ fixed to the mean within-subgraph reference count of papers published
up to 1989. Bounded multistart least squares minimizes the squared error of $\log(1+c)$;
parameter bounds and boundary flags are recorded. Descriptive fits use a random sample of
at most 500 papers published in 1990--2005 with at least 10 citations by age 15, fitted
over ages 1--15; the resulting collapse is in-sample evidence only. For prediction,
papers are selected using only their first five years (at least five citations by age
five; random sample of at most 500), and WSB is compared with a no-ageing preferential
attachment baseline, $c(t)=m[\exp(rt)-1]$, and with linear extrapolation, $c(5)t/5$. The
primary task predicts cumulative citations at age 15, which is observed for every selected
paper; a secondary task predicts the 2020 total with a variable horizon. Neither target
represents lifetime citations. This fixed-horizon design differs from the original study,
which trains on 5 or 10 years and evaluates at age 30 \citep{wang2013quantifying}.
Performance is measured by mean absolute error (MAE), median absolute error and root mean
squared logarithmic error (RMSLE), with paired bootstrap differences (1000 resamples). Fits
at a parameter bound remain in the primary evaluation, and an interior-fit subset serves
as a sensitivity check; a constrained optimum is not regarded as a convergence failure.

\paragraph{P3: citation-target randomization.}
Twenty realizations per model preserve each paper's publication year and its number of
distinct references, which are drawn without replacement from strictly older papers. P3a
selects targets uniformly. P3b weights targets by a nonparametric age kernel fitted by
conditional likelihood to the observed data and removes degree preference. Because P3b is
calibrated on the data, its agreement with the observed ageing is not independent
confirmation. Null means and sample standard deviations are reported; empirical tail
probabilities have a resolution of $1/21$. Both models fix $N(t)$ and $E(t)$ and thus
preserve the densification exponent exactly; they test diameter and target preference but
not densification itself.
